\section{Proof audit and limits of the reduction}
\label{sec:proof-audit}

The seven-rule presentation changes the meaning of an atomic closure.
Consequently its equivalence with the native kernel needs two formula
translations.  It is not an identity-language deletion of one native rule.
This section records the dependency checks that matter for that claim.

\subsection{The order of the arithmetic bootstrap}

Put $A(a,n)=I(a,n;u.Su)$.  The counter is the first argument.
The following order avoids using a derived rule in its own proof.
\begin{enumerate}
\item Computation and equational transport give numerical reflexivity,
      symmetry, transitivity, and term congruence.  For example, if
      $W(t)=I(0,t;u.0)$, the equation $W(t)=t$ is a computation;
      transporting that equation along itself yields $t=t$.
\item A vacuous witness in DI, followed by DE, gives root-directed
      Q-introduction.  This construction uses no numerical induction.
\item Given an equation-induction step and a base proof $E(0)$,
      Q-introduction gives the progressive closure of $E$.  Orbit
      elimination with source $0$ and initializer $t$ gives
      $E(A(0,t))$.  The computation $A(0,t)=t$ yields $E(t)$.
\item Only after this recovery of contextual equation induction do we
      prove $A(a,1)=Sa$.  Orbit elimination with source $a$ and initializer
      $1$ then gives the immediate progress rule
      $\mathsf Q_yE(y),E(a)\vdash E(Sa)$.
\item Full formula transport is proved by structural induction, using
      this immediate progress rule at atomic closure matrices and the
      ordinary base-to-instance elimination at compound matrices.
\end{enumerate}

In particular, the proof does not begin by assuming $I(t,0;u.Su)=t$.
That identity is a consequence of the recovered induction.  Replacing
the target $I(a,n;u.Su)$ in the new elimination rule by
$I(n,a;u.Su)$ would destroy the direct computation in the third item.
The intended rule and the executable representation must retain the
same argument order.

The addition commutation law used to establish soundness in the ordinary
scalar theory is harmless at that stage: the target scalar theory already
has contextual equation induction.  That soundness argument is separate
from the internal recovery of induction in the new calculus.

\subsection{Why equational transport suffices}

For a transport template headed by $\mathsf Q$, the induction hypothesis
applies only to its strictly smaller matrix.  Reverse the supplied
numerical equality to convert the target introduction hypothesis to the
source hypothesis.  Apply the source closure, and convert the resulting
successor instance forwards.  The atomic case uses the immediate progress
rule established above; the compound case uses a zero source in orbit
elimination.  Root-directed Q-introduction closes the argument.

For a template headed by $\mathsf D$, first eliminate the source
$\mathsf D$ at a fresh witness.  Introduce the target $\mathsf D$ at that
same witness, performing the preceding conversion inside its matrix.
Both eigenparameters are chosen outside the transport equality, the
retained context, and the fixed matrix parameters.  Every invocation of
formula transport in these arguments concerns the smaller matrix.
Thus this proof does not silently replace the primitive equational
transport rule by a stronger logical axiom.

The same root invariance is useful for substitution: substituting numerical
terms never turns a literal equation into a compound formula.  The new
introduction and elimination modes therefore survive numerical substitution
without a special witness-zero case.  Renaming eigenparameters first gives
the usual substitution proof by induction on derivations.

\subsection{Root-sensitive closure is not unrestricted congruence}

Padding gives an equation $E$ and a compound formula $P(E)$ that are
interderivable.  Nevertheless one must not infer
\[
               \mathsf Q_yE(y)\ \dashv\vdash\ 
               \mathsf Q_yP(E(y))
\]
merely from that fact.  The first closure is progressive and the second
is conditional.  Numerical transport is legitimate because it preserves
the matrix root.  The binder-congruence lemma used in the formula
composites requires the same mode on its two sides.

The new-side round trip can be justified without assuming logical
completeness for all new formulas in advance.  First translate the old
logical rules for formulas in the image of the old-to-new translation,
allowing an arbitrary passive new context.  Atomic new closures normalize
to image formulas expressing
$\forall y(E(y)\mathbin{\to}E(Sy))$, since $E(y)$ and $E(Sy)$ are
already image formulas.  Atomic witnessed closures then normalize using
the image existential.  Compound cases use the structural induction
hypothesis and same-mode binder congruence.  The translations of compound
formulas remain compound, so that last use of congruence is valid.

One detail is necessary in the atomic witnessed case: elimination of an
image existential must permit an arbitrary new conclusion, not only a
conclusion already known to be an image formula.  Its translated native
macro elimination has precisely this property.  It opens a $\mathsf D$
resource, removes padding and the gate from the witness assumption, and
uses the supplied minor derivation.  The conclusion is only carried
through DE; it is not inspected or introduced beneath a mode-sensitive
binder.  This prevents a circular appeal to the desired normalization
theorem.

\subsection{A concrete obstruction to keeping the old atomic reading}

The change of reading is essential.  A short counterexample uses the
primitive recursive parity term
\[
 p(n)=I(n,0;u.I(u,1;v.0)),\qquad E(y)\equiv p(y)=1.
\]
Direct computations give $p(0)=0$, $p(1)=1$, and $p(2)=0$ in the
standard natural numbers.  Thus
\[
 E(0)\to\forall y E(y)
 \quad\text{is true, whereas}\quad
 \forall y\bigl(E(y)\to E(Sy)\bigr)
 \quad\text{is false}.
\]
Indeed, the old conditional antecedent is false, while the proposed
progress condition fails at $y=1$.

The new orbit elimination would therefore be unsound if its atomic major
kept the old conditional meaning: take source $1$ and initializer $1$.
Its minor $E(1)$ holds and its target is $E(I(1,1;u.Su))=E(2)$, which
fails.  Likewise, merely exporting an ordinary conditional closure from a
progressive witnessed formula discards information needed for a full
formula round trip.  The stronger atomic elimination and the changed
atomic reading must be adopted together.

\subsection{What the comparison establishes}

The new calculus has unrestricted native formula syntax, a single
root-directed introduction, and a direct substitution argument.  These
are structural improvements over the guarded and witness-sensitive
presentations.  They do not by themselves prove a smaller standalone
Python checker.  The orbit elimination has a richer numerical target,
and a faithful implementation must count that cost.  The report's
checker measurements must therefore distinguish this calculus from any
separate compact implementation of an earlier equivalent presentation.

The arguments above are written metatheoretic proofs.  Finite executable
certificates check particular derivations and implementation behavior;
they are not a new Lean or Rocq verification of either formula
interpretation.  No absolute minimality theorem follows from the
seven-rule count or the measured source sizes.
